\chapter{マシンはどう学ぶか}
% full version: chapters/06_dofa.tex

これまでの回路では、結合の強さはすべてエンジニアが選んでいた。生きた脳にエンジニアはいない。結合は、動いているうちに自分で調整されなければならない。しかもマシンが何か役に立つことをするように。これが学習である。

最大の制約は非常に厳しい。各接続部は自分で自分を変え、知ることができるのは身近で起きていることだけだ。送り手が活動しているか、受け手が活動しているか、どんな物質が届いているか。接続部に個別の修正を送ってくれる者はいない。これにより、人工ニューラルネットワークの学習の主要な方法は使えなくなる。そこでは誤差がネットワークを逆向きに伝わり、各結合がそれぞれ自分だけの修正を受け取るからだ。

\section*{一緒に働けば、つながる}

最も単純な規則はこうだ。送り手と受け手が一緒に働けば、その間の結合は強まる。この規則は局所的で、時計を必要としない。少しの忘却を加えれば、この規則は信号の流れの中から最も重要なもの、つまり最も大きく変化するものを見つけ出す。カチッのタイミングを見る版もある。送り手が受け手の\emph{前に}鳴れば、つまり原因でありえたなら結合は強まり、後に鳴れば弱まる。

だがこうした規則には1つの限界がある。学ぶのは\emph{頻繁な}ことであって、\emph{役に立つ}ことではない。隣しか見えない接続部には、行動がジュースをもたらしたのか痛みをもたらしたのかわからない。

\section*{第三の信号}

それを運ぶのがドーパミン、「予想よりよかった」の放送局である。結合が変わるのは、3つのことがそろったときだけにしよう。送り手が活動し、受け手が活動し、そしてドーパミンが届いたとき。

\pencilfig{6}{\pencilart{6}}{2つのニューロンが一緒に働き、第三の信号「予想よりうまくいった」が届いたとき、結合は強まる。}

だが難しい点がある。報酬はすぐには来ない。1、2秒後に来る。その頃には送り手も受け手もとっくに黙っている。接続部には、自分が関わったという記憶が必要だ。そしてそれはある。一緒に働くと、接続部に目印が残る。1、2秒で自然にはがれる付箋のようなものだ。付箋がついている間にドーパミンが来れば、結合が変わる。来なければ、目印は跡形もなく消える。これで宛先の問題が解決する。ドーパミンは全員に共通の1つだが、変わるのは付箋のついた接続部だけだ。

\pencilfig{106}{%
\foreach \y/\d/\t in {0/0.9/{間に合う：結合が強まる},-1.3/3.3/{遅すぎ：変化なし}}{
  \draw[pen] (0,\y) -- (4.6,\y);
  \foreach \s in {0.25,0.33}{\draw[pcBlue, line width=0.8pt] (\s,\y) -- (\s,\y+0.3);}
  \fill[pcAmber, opacity=0.25] (0.35,\y) -- plot[domain=0.35:4.5, samples=30] (\x,{\y+0.8*exp(-(\x-0.35)/0.8)}) -- (4.5,\y) -- cycle;
  \draw[penlite=pcAmber] plot[domain=0.35:4.5, samples=30] (\x,{\y+0.8*exp(-(\x-0.35)/0.8)});
  \draw[pcAmber!80!black, line width=0.9pt] (\d-0.1,\y+0.95) -- (\d+0.07,\y+0.62) -- (\d-0.07,\y+0.6) -- (\d+0.1,\y+0.22);
  \draw[arr=pcAmber!80!black] (\d+0.09,\y+0.27) -- (\d+0.11,\y+0.12);
  \node[lbl, anchor=west] at (4.7,\y+0.3) {\t};}
\node[lbl, text=pcBlue, anchor=south west] at (0.1,0.85) {一緒に活動};
\node[lbl, text=pcAmber!80!black, anchor=south] at (2.3,0.35) {目印が薄れる};
\foreach \x/\l in {0.35/0,1.95/1,3.55/2}{\draw[pcGray, line width=0.5pt] (\x,-1.3) -- (\x,-1.38); \node[lbl, anchor=north] at (\x,-1.38) {\l~s};}
}{一緒に働くと、接続部に1、2秒で消える目印が残る。ドーパミンは目印のある接続部だけを変える。}

\section*{探索としてのノイズ}

この規則はなぜ、でたらめな方向ではなく、よりよい方向へ導くのか。目隠しをして丘の頂上を探しているところを想像してほしい。でたらめに一歩踏み出すと、誰かが「高くなった！」か「低くなった！」と叫ぶ。「高くなった」ならその一歩を確定する。一歩ずつ、一度も地図を見ずに登っていける。脳も同じことをする。ニューロンは少しランダムに揺らぎ、ドーパミンは状況がよくなったかを知らせ、関わった結合は有利な方向へずれる。メッセージの伝達を邪魔していたノイズが、ここでは探索になる。

これで、ドーパミンが報酬ではなく\emph{驚き}を知らせる理由がわかる。うまくいくたびにただ「ジュース！」と叫ぶだけなら、ネットワークがしていたことを、正しいことも間違ったことも、まとめて強化してしまう。私たちのモデルでは、そうしたネットワークは実際に結合を数千倍にまで暴走させ、ときには最悪の選択に永久にはまり込む。期待を差し引くことで、学習は公正になる。

\section*{配線1本の代償}

ブロードキャストによる学習には代償がある。全員に信号は1つで、その中には結果に影響するすべてのニューロンのノイズが混ざっている。ネットワークが大きいほど、各接続部が自分の取り分を取り出すのに時間がかかる。そのためドーパミンが教えるのは、行動を選ぶ比較的小さな回路だけで、巨大な大脳皮質は主に局所的な規則で学習しているらしい。

\section*{驚きを計算するのは誰か}

問題が残る。ドーパミンニューロンは「予想よりよい」をどう計算するのか。必要なのは批評家、つまりこの先どれだけよいことがあるかを常に見積もるニューロン群だ。驚きとは、受け取った報酬に、期待がどれだけ跳ね上がったかを足したものである。電球がつくと期待が跳ね上がり、連打が起きる。約束のジュースが来ると報酬はあるが、同じ瞬間に期待は使い切られるので驚きはない。ジュースが来ないと期待はむなしく崩れ落ち、途切れが生じる。第1章の3つの場合がすべて自然に出てくる。ドーパミンがまさにこう振る舞うことは測定されている。脳がこの驚きを具体的にどう計算しているかは、まだ仮説である。

\fullversion{6}
